The discovery of a new particle consistent with the Standard Model (SM) Higgs boson by the ATLAS~\cite{Aad:2012tfa} and CMS~\cite{Chatrchyan:2012ufa} collaborations
is a major milestone in high-energy physics. However, the underlying nature of electroweak symmetry breaking (EWSB) remains unknown.
Naturalness arguments~\cite{Susskind:1978ms} require that quadratic divergences that arise from radiative corrections to the Higgs boson mass must be canceled by some 
new mechanism in order to avoid fine-tuning. To that effect, several explanations have been proposed in theories beyond the SM (BSM).
In Supersymmetry~\cite{Fayet:1977yc,Djouadi:2005gj,Fayet:1976et}, the cancellation comes from assigning superpartners to the SM bosons and fermions.
A common prediction of supersymmetric models, as in many other BSM scenarios, is the presence of an extended Higgs sector. 
Typical BSM scenarios with an extended Higgs-sector 
explored in the literature include Two-Higgs-Doublets-Models (2HDM)~\cite{Branco:2011iw}, hMSSM~\cite{Djouadi:2015jea,Maiani:2013hud,Djouadi:2013uqa}, and models with Higgs-triplets~\cite{Lee:1973iz,Cheng:1980qt,Schechter:1980gr,Lazarides:1980nt,Mohapatra:1980yp,Magg:1980ut}. 

In 2HDM scenarios, the exclusion limits from searches of BSM Higgs-bosons and SM Higgs coupling measurements~\cite{ATLAS-CONF-2020-053}, restrict the available parameter-space to the so-called alignment limit, in which the SM-like Higgs-boson couplings approach their respective SM values. It has been proposed~\cite{Hajer:2015gka}, that the alignment limit can be probed by searching for BSM heavy Higgs-boson production in association with two top quarks, with the Higgs-boson decaying into a top-quark pair. Additionally, contrarily to inclusive searches for scalars decaying in top-quark pairs, this channel is expected to be insensitive to interference effects with the SM background~\cite{Hajer:2015gka,Gaemers:1984sj}.

Previous searches for heavy neutral Higgs-bosons coupling to top quarks have been performed in Run~II and Run~I proton-proton collisions at center of mass energy of 13, 8 and 7 \tev\ at LHC by the ATLAS~\cite{Aaboud:2017hnm,2023_SSML_4top,2018_Bjets} and CMS~\cite{Sirunyan:2019wxt,Sirunyan_2017} Collaborations; these searches explored a Higgs-boson mass-range between 350~\gev\ and 1.6~\tev\ for several BSM scenarios. Additional strong constraints on extended Higgs sectors are imposed by the measurements of SM Higgs-boson couplings~\cite{Aad:2015pla}. % missing reference in the last word? 
 
This search is also strongly motivated given the recent results obtained from the measurement of the \tttt~production cross-section in the same-sign dilepton and multilepton final state, which was measured to be $24^{+7}_{-6}$~fb for a Standard Model expectation of  $\sigma_{\tttt}=12.0\pm2.4$~fb~\cite{Aad:2020klt}. This search examines different decay channels: opposite-sign dileptons and single lepton final states.   

This note presents a search for new heavy Higgs boson production in a final state involving four-top-quark, $t\bar{t}H/A \to t\bar{t}t\bar{t}$. The search focuses in the mass range of $m_H\in[400,1000]$~\GeV. Above 1 TeV, the effect of t-channel slowly becomes significant above 1 TeV for the interference. It targets a final state with exactly one lepton or two leptons with opposite electric charge and it uses all the data taken  in the full Run 2  with $pp$ collisions at $\sqrt{s} = 13$ TeV using an integrated luminosity of 139 $fb^{-1}$. 
This search follows the analysis strategy developed in Ref.~\cite{Aad:2020klt}, with the extension of the use of new dedicated multivariate discriminant, improved background estimation, 
improved physics object reconstruction, refined model for the description of the systematic uncertainties. 

